% ATTEMPT_STATUS: unresolved
% ATTEMPT_ORIGIN: new research, 2026-10-01
% TOP_PROBLEM: {"id": 1445, "title": "Woodall's Conjecture", "category_id": 3, "category": {"id": 3, "name": "graph_theory", "display_name": "Graph Theory", "description": "Problems involving graphs, networks, and their properties.", "slug": "graph-theory", "order_index": 3, "created_at": "2026-07-31T15:26:25.671Z"}, "status": "open"}
% FIRST_POSTED: 2026-10-01
% Imported from: unsolvedmath_additional_500.tex; UnsolvedMath ID 1445
% !TeX program = lualatex
% Compile twice: lualatex 1445.tex
% Self-contained: no external bibliography, images, or \input files.
% Numeric section labels and visible numbers are original UnsolvedMath IDs.
\documentclass[11pt,a4paper]{article}
\usepackage[margin=24mm,headheight=14pt]{geometry}
\usepackage{fontspec}
\setmainfont{Latin Modern Roman}
\setsansfont{Latin Modern Sans}
\setmonofont{Latin Modern Mono}
\usepackage{amsmath,amssymb,mathtools}
\usepackage{unicode-math}
\setmathfont{Latin Modern Math}
\usepackage{microtype}
\usepackage{enumitem,xurl,needspace}
\usepackage[dvipsnames]{xcolor}
\usepackage{hyperref}
\hypersetup{unicode=true,colorlinks=true,linkcolor=MidnightBlue,urlcolor=MidnightBlue,
 pdftitle={UnsolvedMath Problem 1445},pdfauthor={Source-annotated compilation},bookmarksnumbered=true}
\usepackage{bookmark,tocloft,titlesec,fancyhdr}
\setlength{\cftsecnumwidth}{4.2em}
\cftsetpnumwidth{2.5em}
\cftsetrmarg{4em}
\titleformat{\section}{\Large\bfseries\raggedright}{\thesection}{0.7em}{}
\pagestyle{fancy}\fancyhf{}
\fancyhead[L]{\small\sffamily UnsolvedMath research attempt}
\fancyhead[R]{\small Original catalogue IDs}
\fancyfoot[C]{\thepage}
\setlength{\parindent}{0pt}
\setlength{\parskip}{5pt plus 1pt}
\setlength{\emergencystretch}{3em}
\setcounter{tocdepth}{1}\setcounter{secnumdepth}{1}
\setlist[itemize]{leftmargin=3em,itemsep=4pt,topsep=4pt,parsep=0pt}
\allowdisplaybreaks
\newcommand{\N}{\mathbb{N}}
\newcommand{\Z}{\mathbb{Z}}
\newcommand{\Q}{\mathbb{Q}}
\newcommand{\R}{\mathbb{R}}
\newcommand{\C}{\mathbb{C}}
\newcommand{\F}{\mathbb{F}}
\newcommand{\A}{\mathbb{A}}
\newcommand{\PP}{\mathbb{P}}
\newcommand{\E}{\mathbb{E}}
\newcommand{\eps}{\varepsilon}
\newcommand{\dd}{\,\mathrm d}
\newcommand{\sref}[2]{\hyperlink{source:#1:#2}{[S#2]}}
\newif\ifshowcatalogue
\showcataloguetrue
\newcommand{\cataloguescope}[1]{\ifshowcatalogue
 \paragraph{Statement recorded in the supplied source.}
 #1\fi}
\begin{document}
% UnsolvedMath ID 1445; source catalogue code GRAPH-058

\setcounter{section}{1444}

\section{Woodall’s Dicut--Dijoin Conjecture}\label{1445}

{\small\sffamily UnsolvedMath ID 1445 · Graph Theory}

\subsection{Definitions and mathematical statement}

\paragraph{Shared notation.}Unless a section says otherwise, $\N=\{1,2,\ldots\}$,
$\Z$ is the ring of integers, $\Q,\R,\C$ are the rational, real, and complex
fields, $[N]=\{1,\ldots,\lfloor N\rfloor\}$, and $\log$ is the natural
logarithm. For real functions of a parameter tending to infinity,
$f\ll g$ and $f=O(g)$ mean $|f|\leq Cg$ eventually for some fixed $C$;
$f\gg g$ reverses this comparison; $f=o(g)$ means $f/g\to0$;
$f\sim g$ means $f/g\to1$; and $f\asymp g$ means both order bounds.
A subscript on $O$ or $\ll$ specifies allowed constant dependence.
The expression $n^{o(1)}$ in an upper bound means growth smaller than
$n^\epsilon$ for each fixed $\epsilon>0$. Local definitions govern any
exception, finite-field convention, or use of the same symbol for another
object. An asymptotic claim is not automatically an effective algorithm.



For a vertex set $S$ in a finite digraph, a nonempty directed cut is $\delta^+(S)$ when $\delta^-(S)=\varnothing$. A dijoin is a set of arcs meeting every nonempty directed cut. For a digraph having such a cut, let $\tau$ be its minimum size. The conjecture asks for $\tau$ pairwise arc-disjoint dijoins. No arc weights are imposed. \sref{1445}{1} \sref{1445}{2}

\paragraph{Statement recorded in the supplied source.}
Is the minimum dicut size equal to the maximum number of disjoint dijoins in a directed graph? \sref{1445}{1}

\paragraph{Scope and interpretation.}

The archive’s informal background describes dijoins imprecisely. The cut-hitting definition above is the one in its cited mathematical sources; it is not a prescribed-pair path problem. \sref{1445}{1}

\paragraph{Residual task reported by the archive.}

The embedded review (2026-08-17) records: Close the constant-factor gap and prove or disprove the exact min--max relation. \sref{1445}{1}

\subsection{Short English statement}

Can the arcs of a digraph always contain as many pairwise disjoint directed-cut hitting sets as the size of its smallest directed cut?

\subsection{Research attempt}
\paragraph{Status and scope.}
This attempt leaves Woodall's exact packing conjecture unresolved.
Write $\mathcal D$ for the nonempty directed cuts, assume
$\mathcal D\ne\varnothing$, and put
$\tau=\min_{C\in\mathcal D}|C|$. Internal arcs of strongly connected
components will be irrelevant; parallel arcs are kept as distinct arcs.
The 2025 paper of Cornu\'ejols, Liu and Ravi proves a polynomial-time
packing of at least $\lfloor\tau/6\rfloor$ dijoins. That theorem is an
external benchmark, not derived by the elementary method below.
Its weighted predecessor and the present unweighted conjecture must not
be conflated. \url{https://doi.org/10.1007/s00493-025-00159-x}

\paragraph{The upper bound and the coloring formulation.}
If $J_1,\ldots,J_k$ are disjoint dijoins, each meets a minimum cut $C$.
Their intersections with $C$ are disjoint nonempty sets, so $k\leq|C|=\tau$.
Conversely, color each arc with one of $k$ colors and require that every
cut in $\mathcal D$ contain every color. Each color class is then a
dijoin, and the classes are disjoint. A packing also gives such a
coloring: color the arcs of $J_i$ with color $i$ and give any unused
arcs arbitrary colors. Adding arcs never destroys the property of
being a dijoin. Thus the target is exactly a polychromatic coloring
of the directed-cut family with $\tau$ colors.

This formulation gives a precise finite certificate. Use variables
$x_{a,i}\in\{0,1\}$ with $\sum_{i=1}^{\tau}x_{a,i}=1$ for every arc,
and impose $\sum_{a\in C}x_{a,i}\geq1$ for every $C\in\mathcal D$ and
every color $i$. There may be exponentially many cuts; writing these
constraints is not a polynomial algorithm or a proof of feasibility.

\paragraph{Strong-component contraction.}
If $\delta^-(S)=\varnothing$, no strongly connected component can meet
both $S$ and its complement. Indeed, a path from a vertex outside $S$
to a vertex inside the same component must contain an arc entering $S$.
Therefore $S$ is a union of whole components. Contract the components,
remove internal arcs, and preserve each intercomponent arc individually.
The contracted digraph is acyclic and its nonempty directed cuts are
in bijection with those of the original graph as arc sets. In
particular their minimum size is unchanged. Any coloring of the
contracted arcs meeting all its cuts lifts to the original graph;
internal arcs can be assigned freely. It is consequently sufficient
to solve the problem for acyclic directed multigraphs. The reduction
does not turn a dijoin into a directed path or require prescribed
source-to-sink pairs.

\paragraph{An exact class: a tree of parallel bundles.}
Suppose the underlying simple graph of the contracted digraph is a
tree. For each tree edge $e$, let its $m_e\geq1$ parallel arcs all
point in the same direction. Deleting that tree edge gives two
components; take $S$ to be the component containing its tail. The
corresponding directed cut consists of precisely that bundle, so
$\tau\leq\min_e m_e$. Every nonempty directed cut consists of one or
more whole bundles, and hence has size at least $\min_e m_e$. Thus
\[
 \tau=\min_e m_e.
\]
Color every bundle with all $\tau$ colors, placing at least one arc
of each color in it; this is possible since $m_e\geq\tau$. Every
nonempty directed cut contains a bundle and therefore all colors.
This proves the conjecture for the entire class, with arbitrary tree
orientation and arbitrary positive bundle sizes. The same argument
works componentwise for a forest having at least one bundle. It also
lifts through any strong-component expansion described above.

\paragraph{A probabilistic bound and its exact loss.}
Let $M=|\mathcal D|$. Color the arcs independently and uniformly with
$k\geq2$ colors. For a particular cut $C$ and color $i$, the probability
that $C$ misses $i$ is
$(1-1/k)^{|C|}\leq\exp(-\tau/k)$. There are $Mk$ bad events, so a
coloring and hence $k$ disjoint dijoins exist whenever
\[
 Mk\exp(-\tau/k)<1.
\]
No independence between different bad events is needed for this union
bound. The criterion may yield a packing when $\tau$ is large relative
to $k\log(Mk)$, but it does not reach $k=\tau$: its displayed upper
bound then becomes $M\tau/e$. This failure does not imply a packing
fails to exist. For a single bundle of $\tau$ parallel arcs, assigning
each arc a different color is an explicit optimum even though random
coloring has a coupon-collection obstruction. A sharp argument must
coordinate colors rather than pay separately for every missing color
on every cut. The cited $\lfloor\tau/6\rfloor$ theorem achieves a much
stronger guarantee than this size-dependent union bound.

\paragraph{Verification and residual task.}
The exact companion check enumerates all labeled oriented graphs on
at most four vertices, lists their nonempty directed cuts, and searches
for a $\tau$-coloring satisfying the cut constraints whenever such cuts
exist. This tests the reformulation on every graph in that finite
range. It is not a verification for arbitrary multigraphs, and the
tree-of-bundles theorem instead has the uniform proof given above.
For an arbitrary acyclic digraph, directed cuts can overlap across
many routes and need not contain an entire independently colorable
bundle. A global coloring construction that survives all those
overlaps is the missing step. Status: unresolved.

\subsection{Sources}

{\small

\begin{itemize}

\item[\hypertarget{source:1445:1}{S1}] ULAM.AI, \emph{UnsolvedMath}, supplied \texttt{problems.json}, record 1445 (GRAPH-058), statement and background. The embedded literature-review date is 2026-08-17; that is the archive’s review date, not a fresh certification. Dataset landing page: \url{https://huggingface.co/datasets/ulamai/UnsolvedMath}.

\item[\hypertarget{source:1445:2}{S2}] G. Cornuéjols, S. Liu and R. Ravi, Approximately Packing Dijoins via Nowhere-Zero Flows, Combinatorica 45 (2025), Article 32, \nolinkurl{doi:10.1007/s00493-025-00159-x.} \url{https://doi.org/10.1007/s00493-025-00159-x}. Bibliographic information imported from the supplied record.

\item[\hypertarget{source:1445:3}{S3}] J. Pascal Gollin et al., Disjoint dijoins for classes of dicuts in finite and infinite digraphs, 2022. \url{https://doi.org/10.5070/C62359180}. Bibliographic information imported from the supplied record.

\end{itemize}

}

\label{end:1445}
\end{document}
